% 本文件由 LaTeX 排版 Agent 根据有效 translation.json 自动生成。
% author=Tertullian; work_id=0323; title=致殉道者

\chapter{致殉道者}

% structure: p0001 source=h1 target=omitted reason=page-title-already-rendered-by-parent

Ad Martyras.\par

% structure: p0003 source=h2 target=section reason=relative-heading-rank; base=h2

\section{第一章}

蒙福的候补殉道者，——除了我们的慈母教会从她丰裕的胸怀和各兄弟个人尽力为你们在狱中的肉身需要所作的供应之外，也请接受我提供的一些精神食粮；因为让肉身享受而灵魂挨饿是不好的；不，如果软弱的肉身得到细心照料，那么更软弱的灵魂也应受到关注，这才是公道的。我并非特别有资格劝勉你们；然而，不仅教练和监工，就连外行人，甚至所有随意的人，无需任何必要，也常常习惯从远处用呐喊来激励最娴熟的角斗士，而且有时从单纯的旁观人群中也会产生有用的建议。那么，首先，蒙福的人，不要使那与你们一同进入监狱的圣神忧苦；因为若祂不与你们同去，你们今天就不会在那里。因此，你们要竭尽全力留住祂；让祂带领你们从那里到主那里去。监狱确实是魔鬼的房屋，他在其中安置他的家族。但你们进入它的围墙，正是为了在牠的住处践踏那恶者。你们在外面已经通过激战彻底打败了他；因此，不要让他有任何理由对自己说：\Quote{现在他们在我的领域里；我要用卑鄙的仇恨来试探他们，让他们彼此背叛或分裂。}愿他从你们面前逃走，躲进自己的深渊，萎缩而麻木，就像一条被咒语或烟熏驱赶的蛇。不要让他得逞，使你们在他自己的王国里彼此不合，而要让他发现你们被和谐武装并坚固；因为你们之间的和平就是对他的战斗。有些人在教会中找不到这种和平，便习惯于从被囚的殉道者那里寻求它。因此，你们应当让它与你们同在，并珍视它，守护它，以便你们或许能够将它施予他人。\par

% structure: p0005 source=h2 target=section reason=relative-heading-rank; base=h2

\section{第二章}

其他同样阻碍灵魂的事物，可能一直伴随你们到监狱门口，你们的亲属也可能在那里陪伴你们。从那里和那时起，你们就与世界隔绝了；更不用说与世俗生活的常规及其一切事务隔绝了！不要因这隔绝而惊慌；因为如果我们想到世界才更是真正的监狱，我们就会看到你们是走出监狱，而不是进入监狱。世界有更大的黑暗，蒙蔽人心。世界施加更重的锁链，束缚人的灵魂。世界吐出最污浊的秽气——人的情欲。世界包含更多的罪犯，即全人类。最后，它等待审判，不是总督的审判，而是天主的审判。因此，蒙福的人，你们可以认为自己是从监狱被迁移到了一个安全之所。那里充满黑暗，但你们是光；那里有锁链，但天主使你们自由。那里有难闻的气味，但你们是馨香的芬芳。每天等待审判官，但你们将审判审判官们。那里可能有悲伤，为那叹息世俗享乐的人。监狱外的基督徒已经弃绝了世界，但在监狱中，他也弃绝了监狱本身。你们在世界上身处何处并不重要，因为你们不属于世界。如果你们失去了一些生活的甘甜，那是生意之道：忍受眼前的损失，以便日后获得更大的收益。至此，我还没有提及天主为殉道者准备的赏赐。与此同时，让我们比较一下世界的生活和监狱的生活，看看灵魂在监狱中是否比肉身所失去的获得更多。不，由于教会的关怀和弟兄们的爱，即使肉身在那里也没有失去对其有益的东西，而灵魂还获得了重要的好处。你们没有机会看到异邦的神祇，不会撞见他们的偶像；你们不参与异教徒的节日，即使只是身体上的混杂；你们不被偶像崇拜庆典的污秽烟气所扰；你们不为公开表演的喧嚣、其庆祝者的残暴、疯狂或放荡而痛苦；你们的眼睛不会落在妓院和娼寮上；你们摆脱了冒犯、诱惑和不洁的回忆；你们现在也摆脱了迫害。监狱为基督徒所做的正如旷野为先知所做的。主自己也曾花了很多时间隐居，以便更自由地祈祷，摆脱世界。也是在山上独处时，祂向门徒显示了祂的荣耀。让我们放弃监狱的名称；让我们称它为退隐之所。虽然身体被关住，肉身被束缚，灵魂却通晓一切。因此，灵魂啊，漫游四方；灵魂啊，行走吧，不要将阴凉的小径或长长的柱廊摆在你面前，而要走上那通往天主的道路。每当你的脚步在那里，你就不会在锁链中。当心灵在天上时，腿感觉不到铁链。心灵环绕整个人，并将他带到它愿意去的地方。因为你的财宝在哪里，你的心也必在哪里。\BibleRef{玛 6:21}因此，愿我们的心在那里，我们愿意我们的财宝在那里。\par

% structure: p0007 source=h2 target=section reason=relative-heading-rank; base=h2

\section{第三章}

至福者，现在请赐予恩宠，即使对基督徒而言，监狱亦是令人不悦的；然而，我们正是在对圣事话语的响应中被召叫参加永生天主的战斗。固然，没有士兵带着奢侈品出征，也没有人从舒适的寝室走向战场，而是从轻便狭窄的帐篷出发，在那里必须忍受各种艰苦、粗糙和不适。即使在和平时期，士兵也通过劳苦和艰辛使自己习惯于战争——全副武装行军、在平原上奔跑、挖掘战壕、形成\OriginalTerm{龟甲阵}{testudo}、从事许多艰苦的劳动。一切皆需额上流汗，使身体和心灵不至于因必须从阴凉处走向阳光、从阳光走向严寒、从和平的袍服走向铠甲、从寂静走向喧嚣、从安宁走向骚动而退缩。同样，至福者们，要将你们这命运中的一切艰难视为心志和体力的锻炼。你们即将经历一场崇高的奋斗，在其中，永生天主担任监督，圣神是你们的教练，奖赏是天使本性的永生冠冕、天上的国籍、永恒的荣耀。因此，你们的主耶稣基督，用祂的圣神傅油你们，并引领你们进入竞技场，认为在交战之日前，将你们从本身更舒适的环境中带出，并对你们施以更艰苦的训练是好的，为使你们的力量更强大。因为运动员也被隔离出来接受更严格的训练，以使他们的体力得以增强。他们被禁止奢侈、精美的肉食、更可口的饮品；他们被压制、折磨、耗尽；他们在预备训练中的劳苦越艰巨，胜利的希望就越强大。宗徒说：\BibleQuote{他们是要得可朽坏的花冠。} \BibleRef{格前 9:25} 我们既注目永生的冠冕，便视监狱为我们的训练场，好使在最后审判的终点，我们经多次磨练而被带出，纪律严明；因为德行由艰苦而建立，就如因纵欲享乐而被摧毁。\par

% structure: p0009 source=h2 target=section reason=relative-heading-rank; base=h2

\section{第四章}

从主的话语中我们知道：\BibleQuote{肉躯是软弱的，心神是切愿的。} \BibleRef{玛 26:41} 然而，我们切不可因主承认肉躯的软弱而自欺。正因如此，祂首先宣称心神切愿，为要表明两者中哪一个应当服从另一个——使肉躯服从心神——软弱者服从强者；前者由此从后者获得力量。让心神与肉躯就共同的救恩进行商议，不再思想监狱的磨难，而是思想那作为预备的搏斗与冲突。肉躯或许会惧怕无情的剑、高耸的十字架、野兽的狂怒，以及最可怕的火刑，还有刽子手折磨的一切技巧。但在另一面，让心神向它自己和肉躯清楚地展示：这些事虽然极其痛苦，却已被许多人平静地忍受了——甚至为了名声和荣耀而热切渴望；这不仅是男子如此，女子亦然，为使你们，圣洁的妇女们，堪当你们的性别。若要我逐一列举那些自尽的人，未免太长。至于女子，眼前便有一个著名的例子：被玷污的卢克雷蒂亚，在亲属面前将刀刺入自己怀中，以为自己的贞洁争得荣耀。穆齐乌斯在祭台上焚烧自己的右手，使他的事迹永垂不朽。哲学家们被超越——例如赫拉克利特，他涂上牛粪自焚；恩培多克勒，他跳入埃特纳的火焰中；还有不久前自投火葬堆的佩雷格里努斯。女子甚至也轻视火焰。狄多如此行，以免在深爱的丈夫死后被迫再婚；哈斯德鲁巴的妻子也如此，当迦太基起火时，她不愿看见丈夫匍匐在西庇阿脚下，便带着子女投入故乡被毁的烈火中。雷古鲁斯，一位罗马将军，被迦太基人俘虏，拒绝用大量迦太基俘虏交换，宁愿被交还给敌人。他被塞进一种箱子，四周钉入从外面刺入的钉子，忍受了如此多的十字架酷刑。女子曾自愿寻求野兽，甚至毒蛇，那些比熊或公牛更凶恶的蛇，克娄巴特拉就用它们自尽，以免落入敌人之手。但对死亡的恐惧不如对酷刑的恐惧大。因此，那位雅典的妓女屈服于刽子手，当暴君因她参与阴谋而对她施以酷刑时，她始终没有出卖同谋，最后咬断舌头，吐在暴君脸上，使他确信无论酷刑持续多久都毫无用处。人人都知道至今仍在举行的伟大的拉刻代蒙人的庄严仪式——\OriginalTerm{鞭打}{\GreekText{διαμαστύγωσις}}；在此神圣礼仪中，斯巴达的青少年在祭台前受鞭打，他们的父母和亲属站在旁边，鼓励他们勇敢忍受。因为灵魂而非肉体甘受鞭打，总被认为是更光荣和荣耀的事。但若地上的荣耀——靠心智和体力的活力赢得——如此被看重，以致人们为了同伴的称赞，甚至轻视刀剑、火焰、十字架、野兽、酷刑；那么，为了获得天上的荣耀和神圣的奖赏，这些苦难实在微不足道。若玻璃片如此珍贵，真正的珍珠又价值几何？那么，我们岂不应该以最大的喜乐，为真实的付出别人为虚假的所付出的一切？\par

% structure: p0011 source=h2 target=section reason=relative-heading-rank; base=h2

\section{第五章}

我先放下荣名的动机不谈。所有这些同样残酷痛苦的斗争——你们在世人中间看到的不过是虚荣，实则是某种精神疾病——竟被践踏于脚下。有多少贪图安逸的人，因刀剑的虚妄而丧生？他们真的出于空洞的野心去迎战野兽，还自以为被咬伤和伤疤更迷人。有人卖身赴火，穿着燃烧的束腰外衣跑一段距离；另一些人以最坚韧的肩膀，在猎手的鞭子下行走。主将这些事安置在世上，有福的人们，并非没有理由：那理由是什么？岂不是\Emph{现在}要激励我们，并在那日子定我们的罪——如果我们因惧怕为真理受苦而不敢得救，而别人却出于虚荣热切追求以致丧亡？\par

% structure: p0013 source=h2 target=section reason=relative-heading-rank; base=h2

\section{第六章}

再放下起源于此类坚忍恒心的榜样，让我们转向对人常处境的一份简单默观，或许从那些无论我们愿意与否都会发生、且我们必须决心忍受的事物中，我们可以得到教训。有多少次，烈火吞噬了活人！有多少次，野兽把人撕碎——或是在它们自己的森林里，或是在城市中心，当它们偶然从巢穴逃出时！有多少人倒在强盗的刀下！有多少人先受酷刑，然后又在敌人手中遭受十字架之死，是的，还受尽各种凌辱！人甚至可能为一个人的事业而受苦，却犹豫不肯为天主的事业而受苦。关于这一点，愿现在的时间作证：如此众多的显贵人物仅仅为了一个凡人的事业而丧命，尽管从他们的出身、地位、身体状况和年龄来看，这样的命运似乎最不可能；他们要么因反对他而受他之手，要么因拥护他而受他敌人之手。\par

\SourceRecord{Translated by S. Thelwall.；Ante-Nicene Fathers；Vol. 3.；Edited by Alexander Roberts, James Donaldson, and A. Cleveland Coxe.；Buffalo, NY: Christian Literature Publishing Co.,；1885.；<http://www.newadvent.org/fathers/0323.htm>.}
